% Part: first-order-logic
% Chapter: syntax-and-semantics
% Section: first-order-languages

\documentclass[../../../include/open-logic-section]{subfiles}

\begin{document}

\olfileid{fol}{syn}{fol}

\olsection{پہلے درجے دیاں زباناں}


پہلے درجے دی منطق دیاں لکھتاں اک بنیادی علامتی ذخیرے توں بندیاں نیں،
جیہدے وچ \emph{!!{variable}s}، \emph{!!{constant}s}،
\emph{!!{predicate}s} تے کدی کدی \emph{!!{function}s} ہوندے نیں۔
ایہناں نال منطقی رابطے، مسوّر تے قوساں تے وِرمیاں ورگیاں
رموزی علامتاں ملا کے \emph{ٹرم} تے \emph{!!{formula}s}
بنائے جاندے نیں۔

\begin{explain}
غیر رسمی طور تے !!{predicate}s خاصیتاں تے تعلقاں دے ناں نیں،
!!{constant}s وکھ وکھ شےواں دے ناں نیں، تے !!{function}s
نقشیاں دے ناں نیں۔ !!{identity}~$\eq$ نوں چھڈ کے ایہ
\emph{غیر منطقی علامتاں} نیں تے مل کے اک زبان بناؤندیاں نیں۔
پہلے درجے دی ہر زبان~$\Lang L$ اوہدیاں غیر منطقی
علامتاں توں متعین ہوندی اے۔ سب توں عام صورت وچ
$\Lang L$ وچ ہر قسم دیاں بے انت علامتاں ہوندیاں نیں۔
\end{explain}

عام صورت وچ اسیں پہلے درجے دی منطق وچ ایہ علامتاں ورتدے آں:

\begin{enumerate}
\item منطقی علامتاں
\begin{enumerate}
\item منطقی رابطے:
  \startycommalist
  \iftag{prvNot}{\ycomma $\lnot$ (نفی)}{}%
  \iftag{prvAnd}{\ycomma $\land$ (عطف)}{}%
  \iftag{prvOr}{\ycomma $\lor$ (فصل)}{}%
  \iftag{prvIf}{\ycomma $\lif$ (!!{conditional})}{}%
  \iftag{prvIff}{\ycomma $\liff$ (!!{biconditional})}{}%
  \iftag{prvAll}{\ycomma $\lforall$ (کُلّی مسوّر)}{}%
  \iftag{prvEx}{\ycomma $\lexists$ (وجودی مسوّر)}{}.
\tagitem{prvFalse}{!!{falsity} لئی قضیاتی مستقل~$\lfalse$۔}{}
\tagitem{prvTrue}{!!{truth} لئی قضیاتی مستقل~$\ltrue$۔}{}
\item دو-مقامی !!{identity}~$\eq$۔
\item !!{variable}s دا اک !!{denumerable}s سیٹ: $\Obj v_0$, $\Obj v_1$, $\Obj
  v_2$, \dots
\end{enumerate}
\item غیر منطقی علامتاں، جیہڑیاں پہلے درجے دی منطق دی
\emph{معیاری زبان} بناؤندیاں نیں
\begin{enumerate}
\item $n$-مقامی !!{predicate}s دا ہر $n>0$ لئی اک !!{denumerable}s سیٹ: $\Obj
  A^n_0$, $\Obj A^n_1$, $\Obj A^n_2$, \dots
\item !!{constant}s دا اک !!{denumerable}s سیٹ: $\Obj c_0$, $\Obj c_1$, $\Obj
  c_2$, \dots.
\item $n$-مقامی !!{function}s دا ہر $n>0$ لئی اک !!{denumerable}s سیٹ:
  $\Obj f^n_0$, $\Obj f^n_1$, $\Obj f^n_2$, \dots
\end{enumerate}
\item رموزِ اوقاف: (, )، تے وِرما۔
\end{enumerate}

ساڈیاں زیادہ تر تعریفاں تے نتیجے پہلے درجے دی منطق دی پوری
معیاری زبان لئی بیان کیتے جان گے۔ پر ورتوں دے مطابق اسیں
زبان نوں صرف کجھ !!{predicate}s، !!{constant}s تے
!!{function}s تک محدود وی کر سکدے آں۔

\begin{ex}
حساب دی زبان~$\Lang L_A$ وچ اک دو-مقامی !!{predicate}~$<$،
اک !!{constant}~$\Obj 0$، اک یک-مقامی !!{function}~$\prime$،
تے دو دو-مقامی !!{function}s~$+$ تے~$\times$ ہوندے نیں۔
\end{ex}

\begin{ex}
سیٹ نظریے دی زبان~$\Lang L_Z$ وچ صرف اک دو-مقامی
!!{predicate}~$\in$ ہوندا اے۔
\end{ex}

\begin{ex}
ترتیباں دی زبان~$\Lang L_\le$ وچ صرف اک دو-مقامی
!!{predicate}~$\le$ ہوندا اے۔
\end{ex}

فیر وی ایہ رواجی علامتاں نیں: باقاعدہ طور تے ایہ صرف قائم مقام ناں نیں؛
مثلاً $<$، $\in$ تے $\le$، $\Obj A^2_0$ دے قائم مقام ناں نیں؛
$\Obj 0$، $\Obj c_0$ دا؛ $\prime$، $\Obj f^1_0$ دا؛
$+$، $\Obj f^2_0$ دا؛ تے $\times$، $\Obj f^2_1$ دا۔

\iftag{defNot,defOr,defAnd,defIf,defIff,defTrue,defFalse,defEx,defAll}{%
اُتے متعارف کرائے ابتدائی رابطیاں تے
\iftag{notprvEx,notprvAll}{مسوّر}{مسوّراں} دے نال نال،
اسیں ایہ \emph{تعریف شدہ} علامتاں وی ورتدے آں:
\startycommalist
  \iftag{defNot}{\ycomma $\lnot$ (نفی)}{}%
  \iftag{defAnd}{\ycomma $\land$ (عطف)}{}%
  \iftag{defOr}{\ycomma $\lor$ (فصل)}{}%
  \iftag{defIf}{\ycomma $\lif$ (!!{conditional})}{}%
  \iftag{defIff}{\ycomma $\liff$ (!!{biconditional})}{}%
  \iftag{defAll}{\ycomma $\lforall$ (کُلّی مسوّر)}{}%
  \iftag{defEx}{\ycomma $\lexists$ (وجودی مسوّر)}{}%
  \iftag{defFalse}{\ycomma !!{falsity}~$\lfalse$}{}%
  \iftag{defTrue}{\ycomma !!{truth}~$\ltrue$}}{}.

\begin{tagblock}{defNot,defOr,defAnd,defIf,defIff,defTrue,defFalse,defEx,defAll}
\begin{explain}
تعریف شدہ علامت باقاعدہ طور تے زبان دا حصہ نہیں ہوندی،
بلکہ غیر رسمی مخفف دے طور تے متعارف ہوندی اے: ایہ
اوہناں فارمولیاں نوں مختصر لکھن دی سہولت دیندی اے
جیہڑے صرف ابتدائی علامتاں نال کافی لمّے ہو جاندے۔
ایہ اک صاف فائدہ اے۔ پر اس توں وڈا فائدہ ایہ اے کہ ثبوت
چھوٹے ہو جاندے نیں۔ جے کوئی علامت ابتدائی ہووے، تاں
ثبوتاں وچ اوہدا وکھرا معاملہ کرنا پیندا اے۔ اس لئی
ابتدائی علامتاں جِنیاں ودھ ہون، ساڈے ثبوت اُنّے لمّے ہوندے نیں۔
\end{explain}
\end{tagblock}

% Alternate symbols

\begin{intro}
ممکن اے تہانوں اُتے ورتی اصطلاحات تے علامتاں توں
وکھریاں اصطلاحات تے علامتاں دا پتا ہووے۔ منطق دیاں کتاباں
(تے استاد) ”نفی“ لئی عام طور تے $\sim$، $\neg$ یا~!\
تے ”عطف“ لئی $\wedge$، $\cdot$ یا $\&$ ورتدے نیں۔
”شرطیہ“ یا ”استلزام“ لئی عام علامتاں $\rightarrow$،
$\Rightarrow$ تے $\supset$ نیں۔
\iftag{prvIff,defIff}{”دو طرفہ شرطیہ“، ”دو طرفہ استلزام“
  یا ”(مادی) ہم ارزی“ لئی $\leftrightarrow$،
  $\Leftrightarrow$ تے $\equiv$ علامتاں نیں۔}{}
\iftag{prvFalse,defFalse}{$\lfalse$ دی علامت نوں ”کذب“،
  ”فالسَم“، ”محال“ یا ”تھلّہ“ وی آکھیا جاندا اے۔}{}
\iftag{prvTrue,defTrue}{$\ltrue$ دی علامت نوں ”صدق“،
  ”وِیرَم“ یا ”چوٹی“ وی آکھیا جاندا اے۔}{}

رواج اے کہ !!{constant}s (جیہناں نوں کدی ناں وی آکھدے نیں)
لئی لاطینی حروفِ تہجی دے مُڈھ دے چھوٹے حرف (مثلاً $a$، $b$، $c$)
تے !!{variable}s لئی آخر دے چھوٹے حرف (مثلاً $x$، $y$، $z$)
ورتے جان۔ مسوّر !!{variable}s نال ملدے نیں، مثلاً $x$؛
علامت نگاری دیاں وکھریاں صورتاں وچ کُلّی مسوّر لئی
$\forall x$، $(\forall x)$، $(x)$، $\Pi x$، $\bigwedge_x$
تے وجودی مسوّر لئی $\exists x$، $(\exists
x)$، $(Ex)$، $\Sigma x$، $\bigvee_x$ شامل نیں۔
\end{intro}

\begin{explain}
اسیں سارے قضیاتی عملگر تے دونیں مسوّر زبان دیاں ابتدائی
علامتاں من سکدے آں۔ یا اسیں ابتدائی علامتاں دا چھوٹا
ذخیرہ چُن کے باقی !!{operator}s نوں تعریف شدہ من سکدے آں۔
بولیائی عملگراں دے ”صداقتی-فنکشنی طور تے مکمل“ سیٹاں وچ
$\{ \lnot, \lor \}$، $\{ \lnot, \land \}$ تے $\{ \lnot,
\lif\}$ شامل نیں—ایہناں وچوں ہر سیٹ نوں دوہاں وچوں کسی وی
اک مسوّر نال ملا کے اظہار دے لحاظ نال مکمل پہلے درجے دی
زبان بنائی جا سکدی اے۔

تہانوں دو ہور !!{operator}s دا پتا ہو سکدا اے: شیفر دی لکیر~$|$
(ہنری شیفر دے ناں اُتے)، تے پیئرس دا تیر~$\downarrow$،
جیہنوں کوائن دا خنجر وی آکھدے نیں۔ جدوں انہاں نوں بالترتیب
”نینڈ (nand)“ تے ”نور (nor)“ دے عام معنیاں وچ لیا جائے،
تاں ایہ عملگر اکیلے اکیلے صداقتی-فنکشنی طور تے مکمل نیں۔
\end{explain}

\end{document}
